% Lösungen Binomialverteilung (zusammenbau v0.4, Heft)
\einheitenkopf{Binomialverteilung}

\erg{43}{a) Bernoulli-Kette – zwei Ausgänge, feste Anzahl, gleiches p, unabhängig \quad b) $p = \frac{1}{2} = 0{,}5$ \quad c) $X \sim B(12;\, 0{,}4)$, also $n = 12$ und $p = 0{,}4$}
\erg{44}{a) Zwei Ausgänge je Brief (beschädigt oder nicht); feste Anzahl der ausgewählten Briefe; gleiche Wahrscheinlichkeit für jeden Brief, weil aus sehr vielen Briefen ausgewählt wird und das Nicht-Zurücklegen kaum etwas ändert; die Briefe sind voneinander unabhängig. \quad b) Zwei Ausgänge je Person (Fahrrad oder nicht); feste Anzahl der Befragten; gleiche Wahrscheinlichkeit, weil die Stadt sehr viele Pendler hat und niemand doppelt befragt wird, ohne dass sich der Anteil merklich ändert; die Antworten hängen nicht voneinander ab. \quad c) Nicht binomialverteilt: ohne Zurücklegen aus nur dreißig Laptops ändert sich die Wahrscheinlichkeit für „defekt“ von Zug zu Zug, bei acht von dreißig zu stark, um es zu vernachlässigen; außerdem können höchstens fünf defekte gezogen werden, der Wertebereich reicht nicht bis acht. \quad d) Nicht binomialverteilt: jedes Mitglied kann nur einmal ausgelost werden, also wird ohne Zurücklegen aus einer kleinen Gesamtheit gezogen; die Wahrscheinlichkeit für „Schachspieler“ ändert sich nach jedem Los merklich. \quad e) Erste Aussage richtig: eine Drehung hat nur die zwei Ergebnisse A und B. Zweite Aussage falsch: ein Spiel hat drei Ausgänge – Hauptpreis, Trostpreis oder kein Preis (AB oder BA). \quad f) Erste Aussage richtig: zwei Ausgänge mit fester Wahrscheinlichkeit. Zweite Aussage falsch: eine Runde hat mit genau einem Korb einen dritten Ausgang ohne Preis.}
\erg{45}{a) $P(X = 4)$ – „genau“ verlangt die Einzelwahrscheinlichkeit \quad b) $(6 \text{ über } 2) \cdot 0{,}35^{2} \cdot 0{,}65^{4} \approx 0{,}3280$ \quad c) $\left(\frac{5}{6}\right)^4 = \frac{625}{1\,296} \approx 0{,}4823$}
\erg{46}{a) $0{,}92^{30} \approx 0{,}0820$ \quad b) $0{,}9^{12} \approx 0{,}2824$ \quad c) $\left(\frac{5}{6}\right)^4 + 4 \cdot \frac{1}{6} \cdot \left(\frac{5}{6}\right)^3 = \frac{1\,125}{1\,296} \approx 0{,}8681$ \quad d) $\left(\frac{1}{2}\right)^6 + 6 \cdot \frac{1}{2} \cdot \left(\frac{1}{2}\right)^5 = \frac{7}{64} \approx 0{,}1094$ \quad e) $k = 4$, $n = 7$, $1 - p = \frac{3}{5}$: $P(X = 4) = (7 \text{ über } 4) \cdot \left(\frac{3}{5}\right)^3 \cdot \left(\frac{2}{5}\right)^4$ \quad f) $k = 2$, $n = 8$, $p = 0{,}15$: $P(X = 2) = (8 \text{ über } 2) \cdot 0{,}85^6 \cdot 0{,}15^2$}
\erg{47}{a) $P(X = 6) = P(X \le 6) - P(X \le 5) \approx 0{,}1712$; $P(X < 6) = P(X \le 5) \approx 0{,}4353$ \quad b) $P(X = 11) = P(X \le 11) - P(X \le 10) \approx 0{,}1222$; $P(X < 11) = P(X \le 10) \approx 0{,}3300$ \quad c) $P(X = 51) = (60 \text{ über } 51) \cdot 0{,}85^{51} \cdot 0{,}15^9 \approx 0{,}1429$; höchstens 48 mit Anschluss heißt mindestens 12 ohne: $1 - P(Y \le 11) \approx 0{,}1806$ \quad d) $P(X = 36) = (40 \text{ über } 36) \cdot 0{,}9^{36} \cdot 0{,}1^4 \approx 0{,}2059$; höchstens 33 bestanden heißt mindestens 7 nicht bestanden: $1 - P(Y \le 6) \approx 0{,}0995$}
\erg{48}{a) $P(X \le 4)$ – weniger als fünf heißt höchstens vier \quad b) Grenze 6: $P(X \le 6)$ \quad c) $P(X \le 7) \approx 0{,}3061$ (X: Anzahl, n = 25)}
\erg{49}{a) $P(A) = P(X \le 6) \approx 0{,}3643$; $P(B) = P(X \le 11) - P(X \le 8) \approx 0{,}2674$ \quad b) $P(A) = P(X \le 12) \approx 0{,}2448$; $P(B) = P(X \le 19) - P(X \le 14) \approx 0{,}4395$ \quad c) $60 - H \ge 2H$, also $H \le 20$: $P(H \le 20) \approx 0{,}4516$ \quad d) $X \ge 1{,}5 \cdot (40 - X)$, also $X \ge 24$: $1 - P(X \le 23) \approx 0{,}0405$ \quad e) $0{,}9004$ ist die Wahrscheinlichkeit, dass ein Werkstück als in Ordnung eingestuft wird (beide Pfade); von 50 geprüften Werkstücken werden mit einer Wahrscheinlichkeit von mehr als 60\,\% mindestens 45 als in Ordnung eingestuft. \quad f) $0{,}941$ ist die Wahrscheinlichkeit, dass ein Paket als lesbar erkannt wird (beide Pfade); von 30 Paketen werden mit einer Wahrscheinlichkeit von mehr als 70\,\% mindestens 28 als lesbar erkannt. \quad g) Wie oft darf man höchstens drehen, damit mindestens einmal Gold mit einer Wahrscheinlichkeit von weniger als der Hälfte eintritt? \quad h) Wie viele Lose muss man mindestens kaufen, damit man mit einer Wahrscheinlichkeit von mindestens neunzig Prozent mindestens einen Gewinn hat?}
\erg{50}{a) noch mindestens 23 Treffer in 30 Würfen: $1 - P(Y \le 22) \approx 0{,}2814$ (Y: Treffer im Rest)}
\erg{51}{a) 52\,\% von 400 sind 208, also höchstens 150 Zahl in den restlichen 300 Würfen: $P(Y \le 150) \approx 0{,}5230$ (Y: Zahl im Rest) \quad b) 16\,\% von 1\,200 sind 192, also höchstens 147 Sechsen in den restlichen 900 Würfen: $P(Y \le 147) \approx 0{,}4152$ (Y: Sechsen im Rest)}
\erg{52}{a) n über den Logarithmus – Mindestanzahl bei „mindestens einmal“ \quad b) Nietenwahrscheinlichkeit $0{,}85$: $1 - 0{,}85^n \ge 0{,}9$ \quad c) $n = 4$: mit $0{,}92$ für „fehlerfrei“ ist $P(Y \ge 3) \approx 0{,}7787$ für 3 und $\approx 0{,}9656$ für 4 Gläser (Y: fehlerfreie Gläser)}
\erg{53}{a) höchstens $0{,}0522$ (etwa 5,2\,\%): $(1 - p)^{30} \ge 0{,}2 \Leftrightarrow 1 - p \ge \sqrt[30]{0{,}2}$ \quad b) höchstens $0{,}0172$ (etwa 1,7\,\%): $(1 - p)^{40} \ge 0{,}5 \Leftrightarrow 1 - p \ge \sqrt[40]{0{,}5}$ \quad c) $n = 171$: $P(X > 120) \approx 0{,}8912$ für 170 und $\approx 0{,}9127$ für 171 (X: Anzahl mit Fahrrad) \quad d) $n = 96$: $P(X > 30) \approx 0{,}9434$ für 95 und $\approx 0{,}9516$ für 96 (X: Anzahl der Zugreisenden)}
\erg{54}{a) 26 und 27 ($P(X = 26) \approx 0{,}1001$, $P(X = 27) \approx 0{,}1031$; bei 25 und 28 unter 10\,\%)}
\erg{55}{a) Richtig, es sind 5 Werte: i = 11 bis 15 (etwa 0{,}1021 bis 0{,}1110); die Nachbarn liegen darunter, $P(Y = 10) \approx 0{,}0749$ und $P(Y = 16) \approx 0{,}0867$ \quad b) Falsch, es sind 3 Werte: i = 4, 5, 6 (etwa 0{,}1867, 0{,}1960, 0{,}1633); $P(Y = 3) \approx 0{,}1358$ und $P(Y = 7) \approx 0{,}1108$ liegen darunter}
\erg{56}{a) Fehlschüsse – die 3 steht als Exponent bei der Fehlschusswahrscheinlichkeit \quad b) ausgelassen \quad c) Stelle $4$ (Erwartungswert $16 \cdot 0{,}25 = 4$), Höhe $P(X = 4) \approx 0{,}2252$}
\erg{57}{a) $13$: Erwartungswert 13,5; $P(X = 13) \approx 0{,}1175$ ist größer als $P(X = 14) \approx 0{,}1141$ \quad b) $24$: Erwartungswert 24,5; $P(X = 24) \approx 0{,}0993$ ist größer als $P(X = 25) \approx 0{,}0983$ \quad c) Abbildung 2: das Maximum muss beim Erwartungswert 3 liegen (Abbildung 1 hat es bei 5), und die Säulen müssen zusammen 100 Prozent ergeben (in Abbildung 3 ist allein die Säule bei 3 fast halb so hoch, die Summe liegt weit darüber). \quad d) Abbildung 3: das Maximum muss beim Erwartungswert 9 liegen (Abbildung 2 hat es bei 6), und die Säulen müssen zusammen 100 Prozent ergeben (Abbildung 1 hat das Maximum bei 9, aber die Summe liegt weit über 100 Prozent).}
\erg{58}{a) $P(X = 7) \approx 0{,}85 - 0{,}15 - 0{,}5 = 0{,}20$: Symmetrie um 7,5 gibt $P(X = 6) = P(X = 9)$ und $P(X \ge 8) = 0{,}5$ \quad b) $P(X = 9) \approx 0{,}82 - 0{,}14 - 0{,}5 = 0{,}18$: Symmetrie um 9,5 gibt $P(X = 8) = P(X = 11)$ und $P(X \ge 10) = 0{,}5$ \quad c) $P(Y = 9) \approx 0{,}83 - 0{,}15 - 0{,}5 = 0{,}18$: die Symmetrieachse liegt bei 8,5, also $P(Y = 10) = P(Y = 7)$ und $P(Y \le 8) = 0{,}5$ \quad d) $P(Y = 12) \approx 0{,}80 - 0{,}14 - 0{,}5 = 0{,}16$: die Symmetrieachse liegt bei 11,5, also $P(Y = 13) = P(Y = 10)$ und $P(Y \le 11) = 0{,}5$ \quad e) (I) richtig: bei 35 sind die Säulen von B und C etwa gleich hoch ($\approx 0{,}0463$ und $\approx 0{,}0495$). (II) falsch: C liegt am weitesten rechts, hat also die größte Allergiewahrscheinlichkeit; niedrigere Säulen zeigen nur eine breitere Verteilung. \quad f) (I) richtig: bei 82 sind die Säulen von B und C etwa gleich hoch ($\approx 0{,}0337$ und $\approx 0{,}0308$). (II) falsch: C liegt am weitesten rechts, der Anteil ist dort am größten; die niedrigere Säule zeigt nur eine breitere Verteilung.}
